\begin{lemma}[The Riesz transform is linear on $L^2(\gamma_n)$]
Fix $n \in \mathbb{N}_0$ and let $R_\gamma : L^2(\gamma_n) \to
L^2(\gamma_n;\mathbb{R}^n)$ be the energy-normalized vector Riesz transform of
\noderef{RieszTransform}.  Then $R_\gamma$ is linear over the real scalars, in
the two clauses:
\begin{enumerate}
  \item (additivity) for all $v, w \in L^2(\gamma_n)$,
    $R_\gamma(v + w) = R_\gamma v + R_\gamma w$;
  \item (homogeneity) for every $c \in \mathbb{R}$ and every
    $v \in L^2(\gamma_n)$, $R_\gamma(c\,v) = c\,R_\gamma v$.
\end{enumerate}
Both clauses are unconditional: they hold for every pair of elements and every
real scalar, with no domain restriction and no differentiability assumption on
the arguments.  The scalars are real throughout, $L^2(\gamma_n)$ being a real
Hilbert space, so clause (2) is homogeneity over $\mathbb{R}$ and not merely
over the nonnegative reals.
\end{lemma}
\begin{proof}
By \noderef{RieszTransform}, $R_\gamma$ is the composite
$R_\gamma f = D\,L^{-1/2}\Pi_0 f$: one first forms the potential
$v = L^{-1/2}\Pi_0 f$ and then takes its energy-normalized weak gradient $D$.
So it suffices to see that each of the two factors is linear and that the
intermediate domain conditions are satisfied automatically, after which the
composite of two linear maps is linear.

Step 1 (the first factor is everywhere defined).  The exponent here is
$s = -1/2 \leq 0$, so by the first clause of \noderef{FracPowLinear} every
$v \in L^2(\gamma_n)$ lies in $\operatorname{Dom}(L^{-1/2})$.  Hence no domain
hypothesis has to be carried: $L^{-1/2}\Pi_0$ is defined on all of
$L^2(\gamma_n)$, which is what makes the two clauses above unconditional.

Step 2 (the first factor is linear).  By the remaining two clauses of
\noderef{FracPowLinear}, applied at $s = -1/2$ and using Step 1 to discharge
their domain hypotheses,
\[
  L^{-1/2}\Pi_0 (v + w) = L^{-1/2}\Pi_0 v + L^{-1/2}\Pi_0 w,
  \qquad
  L^{-1/2}\Pi_0 (c\,v) = c\, L^{-1/2}\Pi_0 v
\]
for all $v, w \in L^2(\gamma_n)$ and all $c \in \mathbb{R}$.

Step 3 (the potentials are differentiable).  For every $v \in L^2(\gamma_n)$ the
potential $L^{-1/2}\Pi_0 v$ lies in $\operatorname{Dom}(L^{1/2})$ by
\noderef{PseudoInverseSqrt}, hence in $W^{1,2}(\gamma_n)$ by
\noderef{SqrtDomainSobolev}.  So the second factor $D$ may be applied to every
potential that arises, and again no side condition survives into the statement.

Step 4 (the second factor is linear).  By \noderef{GaussDLinear}, $D$ is
additive and real-homogeneous on $W^{1,2}(\gamma_n)$: for $u, u'$ in that space
and $c \in \mathbb{R}$ one has $D(u + u') = Du + Du'$ and $D(c\,u) = c\,Du$,
the domain hypotheses being supplied by Step 3.

Step 5 (composing).  Fix $v, w \in L^2(\gamma_n)$.  Unfolding the definition of
$R_\gamma$ three times and applying Step 2 and then Step 4,
\[
  R_\gamma(v+w) = D\,L^{-1/2}\Pi_0(v+w)
    = D\bigl(L^{-1/2}\Pi_0 v + L^{-1/2}\Pi_0 w\bigr)
    = D\,L^{-1/2}\Pi_0 v + D\,L^{-1/2}\Pi_0 w
    = R_\gamma v + R_\gamma w ,
\]
which is clause (1).  Likewise, for $c \in \mathbb{R}$ and $v \in L^2(\gamma_n)$,
\[
  R_\gamma(c\,v) = D\,L^{-1/2}\Pi_0(c\,v)
    = D\bigl(c\,L^{-1/2}\Pi_0 v\bigr)
    = c\,D\,L^{-1/2}\Pi_0 v
    = c\,R_\gamma v ,
\]
which is clause (2).  As $v$, $w$ and $c$ were arbitrary, both clauses hold.
\end{proof}
